\section{Experimental Evaluation}
\label{sec:evaluation}

This evaluation does not seek an optimal underwater racing controller. It tests whether the benchmark, official observation contract, referee and team scoring behave as specified. It also reports where the reference baseline fails.

\subsection{Protocol and Metrics}

The evaluation has three experiment groups. First, the HoloOcean clean/current validation covers the three official clean tracks, the two-rover smoke demonstration on Horseshoe Bay and the \texttt{medium} and \texttt{strong} current profiles on Horseshoe Bay. These experiments use the HoloOcean adapter with the BlueROV2 agent, a control timestep of $\Delta t=0.033$\,s, official mode and the unmodified official gate apertures (requirement R4). Because HoloOcean execution is not bit-deterministic, each condition in Table~\ref{tab:validation_results} is evaluated over five seeds.

Second, the deterministic fallback controller/fleet comparison uses the deterministic kinematic adapter on Horseshoe Bay in official mode. It compares two single-rover controllers and three-, four- and five-rover fleets with and without leader-follower coordination. These results are exactly reproducible. Inter-vehicle proximity events on this substrate represent geometric proximity, not contact dynamics.

Third, the HoloOcean diagnostic validation of coordination returns to the HoloOcean adapter. Its scope is only the clean Horseshoe Bay track with three rovers, no currents or obstacles, two start gaps ($8.0$ and $0.0$\,s) and three seeds. It compares the same team with and without leader-follower coordination in diagnostic inter-vehicle mode.

The HoloOcean clean/current metrics are course completion, completed gates, official time, penalized time from \eqref{eq:penalized}, gate/world collisions, out-of-bounds events and stuck events. The deterministic comparison also reports command smoothness, inter-vehicle proximity events and team penalized time. The HoloOcean diagnostic validation of coordination reports team completion, inter-vehicle proximity events, gate/world collisions and team times. Completed-gate and gate/world collision statistics include all seeds. Time statistics include finished runs only and are omitted when no seed finishes.

\subsection{Single-Rover Clean Tracks}
\label{sec:eval_clean}

\begin{table*}[t]
\centering
\caption{Final HoloOcean clean-track, current and homogeneous two-rover fleet results over five seeds per case.}
\label{tab:validation_results}
\mratablestyle
\begin{tabular}{llrlrrr}
\toprule
Setting & Controller and track & Seeds & Status & Gates & Time (s) & Collisions \\
\midrule
\multirow{6}{*}{Clean}
 & Continuous servo, Horseshoe Bay & 5 & 5/5 FIN & 12.0/12 & $194.5\pm4.3$ & 0.0 \\
 & Center-then-commit, Horseshoe Bay & 5 & 5/5 FIN & 12.0/12 & $197.7\pm6.3$ & 0.0 \\
 & Continuous servo, Vertical Serpent & 5 & 5/5 FIN & 17.0/17 & $236.9\pm5.7$ & 0.0 \\
 & Center-then-commit, Vertical Serpent & 5 & 4/5 FIN & 14.8/17 & $241.1\pm1.6$ & 0.0 \\
 & Continuous servo, Mixed Endurance & 5 & 2/5 FIN & 17.0/22 & $394.7\pm2.8$ & 0.0 \\
 & Center-then-commit, Mixed Endurance & 5 & 5/5 FIN & 22.0/22 & $410.7\pm7.9$ & 0.0 \\
\midrule
\multirow{4}{*}{Currents}
 & Continuous servo, Horseshoe medium & 5 & 2/5 FIN & 8.4/12 & $337.9\pm13.6$ & 43.4 \\
 & Center-then-commit, Horseshoe medium & 5 & 3/5 FIN & 8.8/12 & $217.3\pm1.2$ & 25.2 \\
 & Continuous servo, Horseshoe strong & 5 & 0/5 FIN & 3.0/12 & \textemdash & 0.0 \\
 & Center-then-commit, Horseshoe strong & 5 & 0/5 FIN & 3.2/12 & \textemdash & 0.8 \\
\midrule
\multirow{2}{*}{Fleet, gap 90 s}
 & Continuous servo, team aggregate & 5 & 5/5 FIN & 24.0/24 & $303.1\pm2.9$ & 0.0 \\
 & Center-then-commit, team aggregate & 5 & 5/5 FIN & 24.0/24 & $308.0\pm5.2$ & 0.0 \\
\bottomrule
\end{tabular}
\mratablenote Times are mean$\pm$population standard deviation over finished runs. Gates and collisions are five-seed means; fleet values are team-level and collisions combine gate and world events. Medium-current penalized times are $692.9\pm53.6$\,s for continuous servo and $223.9\pm1.5$\,s for center-then-commit.
\end{table*}


Table~\ref{tab:validation_results} reports the clean-track results over five seeds. The official rule baseline completes all three courses in every seed: $12/12$, $17/17$ and $22/22$ gates, or 100\% completion. No run records an out-of-bounds or stuck event. Mean gate/world collisions remain below one per run: $0.6\pm0.5$, $0.2\pm0.4$ and $0.8\pm0.8$ on Horseshoe Bay, Vertical Serpent and Mixed Endurance. Official times are stable across seeds at $232.4\pm7.5$\,s, $259.7\pm27.8$\,s and $540.6\pm7.7$\,s. These results establish that the official clean tasks are feasible and reproducible under the official observation contract.

\subsection{Staggered Fleet and Team Scoring}

The bottom block of Table~\ref{tab:validation_results} reports the stable two-rover staggered demonstration on Horseshoe Bay over five seeds. Both rovers finish $12/12$ gates in every seed. The first rover finishes in $230.7\pm5.4$\,s with $1.4\pm0.5$ gate/world collisions. The second finishes in $236.5\pm1.4$\,s with $1.2\pm0.4$ gate/world collisions. The team completes $24/24$ gates in every seed. Its elapsed time is $340.7\pm1.4$\,s from first release to last finish, and its penalized time is $353.7\pm3.2$\,s. No inter-vehicle proximity events are recorded because the $90$\,s start gap deliberately separates the rovers. The run therefore tests multi-rover spawning, staggered release, independent per-rover referee state, ranking and team aggregation without physical interaction. It is a low-contact smoke demonstration of the fleet and team machinery.

\subsection{Behavior Under Currents}
\label{sec:eval_currents}

Clean-track performance degrades as current strength increases. On Horseshoe Bay, the rule baseline finishes only $3/5$ seeds under the \texttt{medium} profile. Across all five seeds, it completes $10.8\pm1.5$ of $12$ gates and accumulates $82.6\pm74.1$ gate/world collisions. Among the three finished seeds, official time is $315.0\pm42.3$\,s and penalized time is $428.3\pm98.5$\,s. Under the \texttt{strong} profile, the baseline never finishes. Every seed reaches only $3/12$ gates within the time limit and accumulates $11.0\pm1.1$ gate/world collisions. These results quantify the degradation under the official observation contract with the gates, margins and current profiles unchanged (requirement R4). Designing controllers that reject the current remains future work (Section~\ref{sec:discussion}).

\subsection{Controller Comparison and Fleet Coordination}
\label{sec:eval_comparison}

This experiment asks whether the benchmark separates controllers and whether coordination controls fleet interaction. It uses the engine-free deterministic kinematic adapter on Horseshoe Bay in official mode. The official gate apertures, referee rules and inter-vehicle thresholds remain unchanged (R4). The results are exactly reproducible and isolate the control logic. Inter-vehicle proximity events represent geometric proximity on this substrate, not HoloOcean contact dynamics.

The two controllers occupy different operating points. Both finish all $12$ gates. \cmd{acoustic\_baseline} takes $120.1$\,s, whereas the conservative \cmd{smooth\_gate\_baseline} takes $198.0$\,s, or $1.6\times$ as long. The mean per-step change in surge and yaw commands is used as a lower-is-smoother proxy. It drops from $0.0159$ to $0.0035$, a reduction of roughly $4.5\times$. The benchmark therefore distinguishes the two legal controllers by completion time and behavior.

The fleet comparison uses a slower \cmd{smooth\_gate\_baseline} leader with faster \cmd{acoustic\_baseline} followers. The start gap is $8$\,s, lateral spacing is $1.5$\,m and scoring uses \cmd{penalize} mode. Without coordination, the followers overtake the leader and trigger the proximity detector. Teams of three, four and five record $2$, $3$ and $4$ inter-vehicle proximity events, respectively (Table~\ref{tab:coordination}). Leader-follower coordination removes all of these events, matching the single-rover result, while every rover finishes the full gate sequence. Team time increases because the followers pace behind the leader instead of overtaking it. This is the measured safety-throughput tradeoff.

\begin{table}[t]
\centering
\caption{Fleet coordination on the deterministic kinematic substrate (Horseshoe Bay, official mode, penalize mode for inter-vehicle proximity events). A staggered heterogeneous team has a slower \texttt{smooth\_gate\_baseline} leader and faster \texttt{acoustic\_baseline} followers. Each team size is run with and without leader-follower coordination. Every rover finishes in every condition.}
\label{tab:coordination}
\footnotesize
\setlength{\tabcolsep}{5pt}
\renewcommand{\arraystretch}{1.15}
\begin{tabular}{rlrr}
\toprule
Team & Condition & Inter-vehicle $\downarrow$ & Team penalized (s) \\
\midrule
\multirow{2}{*}{3} & no coordination & 2 & 216.1 \\
 & leader-follower & \textbf{0} & 251.1 \\
\midrule
\multirow{2}{*}{4} & no coordination & 3 & 221.9 \\
 & leader-follower & \textbf{0} & 274.1 \\
\midrule
\multirow{2}{*}{5} & no coordination & 4 & 226.8 \\
 & leader-follower & \textbf{0} & 295.5 \\
\bottomrule
\end{tabular}
\end{table}


The two-gate yield margin of \eqref{eq:yield} is necessary. With a one-gate margin ($\Delta g=1$), a follower can reach a gate just as its predecessor clears it, so a full gate of physical spacing is not preserved. The four-rover team then records $6$ inter-vehicle proximity events, twice the uncoordinated team's $3$. With $\Delta g=2$, it records none. The coordination also tolerates channel loss. Repeating the four-rover coordinated run with $20\%$ per-link packet loss still yields no inter-vehicle proximity events and every rover finishes. Periodic heartbeats and the freshness window absorb the dropped packets. These results support the separation argument of Section~\ref{sec:coordination}: a two-gate progress lead becomes physical spacing. The deterministic kinematic substrate isolates this proximity logic but does not model contact. The next subsection tests the same policy in HoloOcean.

\subsection{HoloOcean Diagnostic Validation}
\label{sec:eval_holoocean_coord}

HoloOcean diagnostic validation tests the same coordination under contact dynamics. It uses only the clean Horseshoe Bay track in official mode with three rovers, no currents or obstacles, two start gaps ($8.0$ and $0.0$\,s) and three seeds. The team has a \cmd{smooth\_gate\_baseline} leader and \cmd{acoustic\_baseline} followers. Inter-vehicle mode is diagnostic and no parameters are retuned. All six coordinated runs finish with $36/36$ team gates. They record no inter-vehicle proximity events, gate/world collisions, out-of-bounds events or stuck events. Penalized time ranges from $130.7$--$131.3$\,s (Table~\ref{tab:holoocean_coordination}). Without coordination, faster followers crowd the slower leader. These runs record one or two inter-vehicle proximity events and up to $531$ gate/world collisions. In two of the six runs, both with the $8$\,s gap, the leader remains unfinished at $4/12$ gates.

Contact dynamics change the tradeoff observed on the deterministic kinematic substrate. There, coordination removes inter-vehicle proximity events at a small time cost. In HoloOcean, the same crowding produces costly gate/world collisions, so coordination becomes a net gain. Finished uncoordinated runs are slightly faster in raw time ($104$--$118$\,s versus $\sim\!131$\,s when coordinated). Once gate/world collisions are scored, their penalized times rise to $222$--$404$\,s. Two of the six uncoordinated runs do not finish. This HoloOcean diagnostic validation is deliberately limited to one clean track, three rovers, two start gaps and three seeds. Extending it to other tracks, currents, obstacles and larger fleets remains future work (Section~\ref{sec:discussion}).

\begin{table*}[t]
\centering
\caption{Three-vehicle coordination on clean Horseshoe Bay over three seeds per condition.}
\label{tab:holoocean_coordination}
\mratablestyle
\begin{tabular}{llrlrrr}
\toprule
Setting & Policy & Seeds & Status & Gates $\uparrow$ & Team time (s) $\downarrow$ & Events $\downarrow$ \\
\midrule
\multirow{3}{*}{Gap 0 s}
 & No coordination & 3 & 3/3 FIN & 36.0/36 & $219.7\pm2.5$ & 9.3 / 2.3 / 0.0 \\
 & LF(1) & 3 & 3/3 FIN & 36.0/36 & $262.3\pm7.7$ & 0.0 / 0.0 / 0.0 \\
 & LF(2) & 3 & 3/3 FIN & 36.0/36 & $288.3\pm1.3$ & 0.0 / 0.0 / 1.0 \\
\midrule
\multirow{3}{*}{Gap 8 s}
 & No coordination & 3 & 2/3 FIN & 34.3/36 & $254.0\pm24.5$ & 127.3 / 2.0 / 0.0 \\
 & LF(1) & 3 & 3/3 FIN & 36.0/36 & $266.0\pm8.3$ & 0.0 / 0.0 / 0.0 \\
 & LF(2) & 3 & 3/3 FIN & 36.0/36 & $285.7\pm3.8$ & 0.0 / 0.0 / 0.0 \\
\bottomrule
\end{tabular}
\mratablenote Times are raw team elapsed mean$\pm$sample standard deviation over full-team finishes; gates and events are three-seed means. Events are GW / IV / S: gate and world collisions, inter-vehicle proximity and stuck. No run records an out-of-bounds event. Penalized time differs for no coordination at gaps 0 and 8 s ($266.4\pm64.1$ and $564.0\pm463.0$\,s) and LF(2) at gap 0 s ($303.3\pm1.3$\,s).
\end{table*}


\subsection{Reproducibility}

Every run emits a structured event log and summary. The repository provides the exact commands used here~\cite{holodronecompetition2026}. A single configuration-driven entry point and per-track JSON files make each official run reproducible from one command. The fleet demonstration and unit-test suite are also scripted, so Table~\ref{tab:validation_results} can be regenerated end to end. One engine-free command produces the deterministic comparison in Table~\ref{tab:coordination} bit-for-bit.
